\chapter{Elastic Bodies And Gauge Kinematics Of Deformation}
\label{chap:elastic-gauge}

\section{Why This Chapter Is Separate From Fluids}
\label{sec:elastic-vs-fluid}

Fluids and elastic bodies share continuum balance laws, but they organize
material response differently. A fluid has no preferred static shear shape: its
deviatoric stress is tied to the rate of deformation. An elastic body stores
energy in deformation itself. This chapter therefore returns to the material
map
\[
  \varphi:B\times I\to\R^3,
  \qquad
  x=\varphi(X,t),
\]
and treats \(X\in B\) as a persistent label of a material point.

The second purpose of the chapter is to introduce the gauge-theoretic
description of deforming bodies in the spirit of Shapere and Wilczek.\footnote{
A. Shapere and F. Wilczek, ``Gauge kinematics of deformable bodies,''
American Journal of Physics 57, 514--518 (1989),
\href{https://doi.org/10.1119/1.15986}{doi:10.1119/1.15986}. Their closely
related low-Reynolds-number swimming formulation appears in
\href{https://doi.org/10.1017/S002211208900025X}{doi:10.1017/S002211208900025X}.}
The central idea is simple and powerful: a deformable body has a space of
shapes, and choosing axes for each shape is a gauge choice. Cyclic changes of
shape can produce a net rotation or translation, not because of a force impulse,
but because a connection over shape space has nonzero curvature.

\begin{table}[htbp]
\centering
\caption{Elasticity, Cosserat, and Shapere-Wilczek gauge notation.}
\label{tab:elastic-gauge-notation}
\small
\begin{tabular}{p{0.18\textwidth}p{0.48\textwidth}p{0.25\textwidth}}
\toprule
Symbol & Meaning & Type\\
\midrule
\(B\) & reference body & material domain\\
\(\partial B\) & boundary of the reference body & material surface\\
\(X^A\) & material coordinate & coordinate on \(B\)\\
\(x^i\) & spatial coordinate & coordinate in \(\R^3\)\\
\(\varphi(X,t)\) & deformation map & \(B\to\R^3\)\\
\(F^i{}_A\) & deformation gradient & \(\partial_A\varphi^i\)\\
\(J\) & local volume ratio & \(\det F\)\\
\(C_{AB}\) & right Cauchy-Green tensor & \(F^i{}_AF^i{}_B\)\\
\(E_{AB}\) & Green strain & \(\frac12(C_{AB}-\delta_{AB})\)\\
\(\xi(X,t)\) & small displacement field & \(\varphi=X+\xi\)\\
\(e_{AB}\) & infinitesimal strain & symmetric gradient of \(\xi\)\\
\(W(F,X)\) & stored energy density & energy per reference volume\\
\(\rho_0(X)\) & reference mass density & mass per reference volume\\
\(P^A{}_i\) & first Piola-Kirchhoff stress & \(\partial W/\partial F^i{}_A\)\\
\(\sigma_{ij}\) & Cauchy stress & spatial stress tensor\\
\(\lambda,\mu\) & Lame moduli & elastic constants\\
\(c_P,c_S\) & longitudinal and shear wave speeds & speed\\
\(g,s,\mathcal S\) & rigid placement, shape, and shape space & gauge kinematics\\
\(\xi=g^{-1}\dot g\) & rigid velocity in moving frame & Lie-algebra element\\
\(M(s),m(s)\) & rigid and shape blocks of kinetic metric & matrices/operators\\
\(\Pi\) & momentum conjugate to rigid motion & element of \(\mathfrak g^*\)\\
\(\mathcal A\) & mechanical/gauge connection & Lie-algebra-valued one-form\\
\(\mathcal F\) & curvature of \(\mathcal A\) & Lie-algebra-valued two-form\\
\(\Gamma\) & closed loop in shape space & stroke/cycle\\
\(\mathcal P\) & path-ordering symbol & ordered exponential\\
\(R(X,t)\) & local material frame & \(\SO(3)\)-valued field\\
\(A_A,A_t\) & spatial and time connection components & \(\mathfrak{so}(3)\)-valued\\
\(\Gamma_A\) & translational strain in local frame & \(R^{-1}\partial_A\varphi\)\\
\(T^a\) & torsion two-form & dislocation density model\\
\(\kappa,B,\kappa_0\) & rod curvature, bending stiffness, intrinsic curvature & rod variables\\
\bottomrule
\end{tabular}
\end{table}

\paragraph{Roadmap.}
The chapter has two intertwined narratives. The first is classical elasticity:
deformation gradients measure local stretching, compatibility distinguishes
true displacement fields from incompatible strain data, and a hyperelastic
action gives elastodynamic equations and stress. The second is gauge
kinematics: once a body can change shape, choosing body axes becomes a
shape-dependent convention, and cyclic shape change can produce rigid motion
through holonomy. The Cosserat and defect sections show how these two
narratives meet in local frame fields.

The guiding distinction is between intrinsic deformation and frame description.
Quantities such as \(C=F^TF\) and the Green strain do not change under a rigid
rotation of the observer. Quantities such as local frame components and
connection coefficients do change, but in a controlled gauge-covariant way. The
goal is not to avoid frames; it is to know exactly which statements depend on
them.

\paragraph{Vocabulary bridge.}
In the elasticity sections, ``configuration'' means a deformation map
\(\varphi:B\to\R^3\). In the gauge sections, the same physical configuration is
split into a rigid placement \(g\) and a shape \(s\). This split is useful, but
not unique: changing the convention for body axes changes \(g\) and changes the
connection coefficients \(\mathcal A_\alpha\). The observable statement is the
reconstructed rigid motion around a closed shape cycle, especially its
curvature-controlled small-loop limit. Thus the gauge language is not an extra
force law; it is a precise bookkeeping system for frame-dependent descriptions
of deformation.

\section{Deformation Gradient And Strain}
\label{sec:deformation-gradient-strain}

Elasticity begins by comparing nearby material points before and after
deformation. The object that does this comparison is not the displacement
itself but its derivative: two deformations that differ by a rigid translation
have the same local strain. This is why the deformation gradient \(F\), and the
metric-like tensor \(C=F^TF\), are the natural first objects.

The deformation gradient is the local linearization of the deformation map:
\[
  F^i{}_A=\frac{\partial \varphi^i}{\partial X^A}.
\]
If \(\delta X\) is a small material vector, then its spatial image is
\[
  \delta x^i=F^i{}_A\delta X^A.
\]
The right Cauchy-Green tensor measures squared lengths after deformation:
\[
  |\delta x|^2
  =
  C_{AB}\delta X^A\delta X^B,
  \qquad
  C_{AB}=F^i{}_AF^i{}_B.
\]
The Green strain is
\[
  E_{AB}=\frac12(C_{AB}-\delta_{AB}).
\]
It vanishes for every rigid motion
\[
  \varphi(X)=a+QX,\qquad Q\in\SO(3),
\]
because \(F=Q\) and \(Q^TQ=I\). This invariance is why \(E\), not \(F-I\), is a
geometric strain measure for finite deformations.

For small displacement
\[
  \varphi(X,t)=X+\xi(X,t),
\]
one has
\[
  F^i{}_A=\delta^i_A+\partial_A\xi^i.
\]
Keeping only first-order displacement gradients,
\[
  E_{AB}
  =
  \frac12(\partial_A\xi_B+\partial_B\xi_A)
  +
  O(|\nabla\xi|^2).
\]
The infinitesimal strain is therefore
\[
  \boxed{
  e_{AB}=\frac12(\partial_A\xi_B+\partial_B\xi_A).
  }
\]

\section{Compatibility}
\label{sec:elastic-compatibility}

Not every matrix field is a deformation gradient. If
\[
  F^i{}_A=\partial_A\varphi^i
\]
for a smooth deformation, then mixed partial derivatives commute:
\[
  \boxed{
  \partial_AF^i{}_B-\partial_BF^i{}_A=0.
  }
\]
On a simply connected body, this curl-free condition is also enough to recover
\(\varphi\) from \(F\), up to a rigid translation.

For small strain, the compatibility question is subtler because \(e\) is only
the symmetric part of \(\nabla\xi\). Given a symmetric tensor field \(e_{ij}\),
we ask whether there exists a displacement \(\xi\) such that
\[
  e_{ij}=\frac12(\partial_i\xi_j+\partial_j\xi_i).
\]
The answer is controlled by the Saint-Venant tensor:
\[
  \boxed{
  \epsilon_{ik\ell}\epsilon_{jmn}
  \partial_k\partial_m e_{\ell n}=0.
  }
\]
This is often written
\[
  \operatorname{curl}\operatorname{curl} e=0.
\]
When the condition fails, the local strain field cannot be assembled into a
single-valued displacement without cuts, defects, residual stress, or a
non-Euclidean reference geometry.

\begin{remark}[Geometric meaning]
The compatibility equations are the linearized version of flatness. A strain
field defines a metric \(g_{ij}=\delta_{ij}+2e_{ij}\). If this metric is the
pullback of the Euclidean metric by an embedding, its Riemann curvature must
vanish. Linearizing that curvature gives Saint-Venant compatibility.
\end{remark}

\section{Hyperelastic Action And Piola Stress}
\label{sec:hyperelastic-action}

A hyperelastic material has stored energy density \(W(F,X)\) per reference
volume. Let \(\rho_0(X)\) be reference mass density. The action is
\[
  S[\varphi]
  =
  \int_{t_0}^{t_1}\int_B
  \left(
    \frac12\rho_0|\partial_t\varphi|^2-W(F,X)
  \right)\dd^3X\,\dd t.
\]

\begin{assumption}[Hyperelasticity]
The stress is derived from stored energy:
\[
  P^A{}_i=\frac{\partial W}{\partial F^i{}_A}.
\]
This excludes plasticity, viscosity, fracture, and history dependence. Those
effects require additional internal variables or dissipation laws.
\end{assumption}

\begin{derivation}[Elastodynamic equation]
Let \(\delta\varphi=\eta\), with \(\eta=0\) at the time endpoints. Since
\[
  \delta F^i{}_A=\partial_A\eta^i,
\]
the first variation is
\[
\begin{aligned}
  \delta S
  &=
  \int_{t_0}^{t_1}\int_B
  \left(
    \rho_0\dot\varphi_i\dot\eta_i
    -
    P^A{}_i\partial_A\eta^i
  \right)\dd^3X\,\dd t\\
  &=
  \int_{t_0}^{t_1}\int_B
  \left(
    -\rho_0\ddot\varphi_i
    +
    \partial_A P^A{}_i
  \right)\eta^i\dd^3X\,\dd t
  -
  \int_{t_0}^{t_1}\int_{\partial B}
  P^A{}_iN_A\eta^i\,\dd S\,\dd t.
\end{aligned}
\]
For arbitrary interior variations,
\[
  \boxed{
  \rho_0\ddot\varphi_i=\partial_A P^A{}_i.
  }
\]
If a boundary traction \(t_i\) per reference area is prescribed, the natural
boundary condition is
\[
  P^A{}_iN_A=t_i.
\]
\end{derivation}

\section{Spatial Stress And Frame Indifference}
\label{sec:spatial-stress-frame-indifference}

The Cauchy stress \(\sigma\) is the physical stress acting across spatial
surfaces. It is related to Piola stress by
\[
  \boxed{
  \sigma_{ij}
  =
  \frac1J P^A{}_iF^j{}_A.
  }
\]
This relation is the Piola transform: it converts force per reference area to
force per current spatial area.

A stored energy is frame indifferent if a superposed rigid rotation does not
change it:
\[
  W(QF,X)=W(F,X)\qquad\text{for all }Q\in\SO(3).
\]
Let \(Q(\eps)=I+\eps A+O(\eps^2)\), with \(A^T=-A\). Differentiating
frame indifference at \(\eps=0\) gives
\[
  0=
  P^A{}_iA^i{}_jF^j{}_A.
\]
Since \(A\) is any skew matrix, the tensor \(PF^T\) must be symmetric:
\[
  P^A{}_iF^j{}_A=P^A{}_jF^i{}_A.
\]
Therefore \(\sigma=\sigma^T\). Thus symmetry of stress follows from rotational
invariance of stored energy, matching angular-momentum balance.

\section{Linear Isotropic Elasticity}
\label{sec:linear-isotropic-elasticity}

For small strain in an isotropic homogeneous solid, the quadratic energy density
is
\[
  W(e)=\frac{\lambda}{2}(\tr e)^2+\mu e_{ij}e_{ij},
\]
where \(\lambda,\mu\) are Lame moduli. The stress is
\[
  \sigma_{ij}
  =
  \frac{\partial W}{\partial e_{ij}}
  =
  \lambda(\tr e)\delta_{ij}+2\mu e_{ij}.
\]
The linearized equation of motion is
\[
  \rho_0\partial_t^2\xi_i=\partial_j\sigma_{ij}+\rho_0 b_i.
\]
Substituting the stress-strain relation,
\[
\begin{aligned}
  \partial_j\sigma_{ij}
  &=
  \lambda\partial_i(\partial_k\xi_k)
  +
  \mu\partial_j(\partial_i\xi_j+\partial_j\xi_i)\\
  &=
  (\lambda+\mu)\partial_i(\nabla\cdot\xi)
  +
  \mu\Delta\xi_i.
\end{aligned}
\]
Hence
\[
  \boxed{
  \rho_0\partial_t^2\xi
  =
  (\lambda+\mu)\nabla(\nabla\cdot\xi)
  +
  \mu\Delta\xi
  +
  \rho_0 b.
  }
\]
Using
\[
  \Delta\xi=\nabla(\nabla\cdot\xi)-\nabla\times\nabla\times\xi,
\]
one sees two wave families:
\[
  c_P=\sqrt{\frac{\lambda+2\mu}{\rho_0}},
  \qquad
  c_S=\sqrt{\frac{\mu}{\rho_0}}.
\]
Longitudinal waves change volume; shear waves change shape at nearly fixed
volume.

\section{Gauge Kinematics Of Deformable Bodies}
\label{sec:gauge-kinematics-deformable-bodies}

The preceding sections describe how an elastic body deforms relative to a
reference body. Gauge kinematics asks a different but compatible question: if
the body changes shape internally, how much of its observed motion is a genuine
rigid displacement, and how much is a convention for choosing body axes at each
shape? This is the same distinction that appears in gauge theory: local
descriptions depend on a frame choice, while holonomy around a closed loop is
observable.

The Wilczek-Shapere point of view begins with a finite deformable body moving in
space. Separate its configuration into:
\[
  g\in G,\qquad s\in\mathcal S,
\]
where \(G\) is the rigid-motion group, usually \(\SE(3)\) or \(\SO(3)\), and
\(\mathcal S\) is shape space. A point of the full configuration space is a
rigid placement \(g\) together with an internal shape \(s\).

The subtlety is that the split into \(g\) and \(s\) is not unique. To say what
the body's orientation is at a given shape, one must choose body axes for that
shape. A shape-dependent change of axes
\[
  h:\mathcal S\to G
\]
changes the representative \(g\). This is exactly a gauge choice.

\begin{figure}[htbp]
\centering
\begin{tikzpicture}[scale=1.0, >=Latex]
  \draw[->] (-0.2,0) -- (6.5,0) node[right] {\(s^1\)};
  \draw[->] (0,-0.2) -- (0,4.1) node[above] {\(s^2\)};
  \draw[thick, blue, ->] (1.2,1.2) .. controls (2.8,3.4) and (5.0,3.1) .. (5.0,1.4);
  \draw[thick, blue, ->] (5.0,1.4) .. controls (3.9,0.4) and (2.2,0.3) .. (1.2,1.2);
  \node[blue] at (3.3,3.55) {shape cycle};
  \begin{scope}[shift={(1.2,1.2)}, scale=0.28]
    \draw[fill=gray!20, thick] (-0.8,0) ellipse (0.45 and 0.25);
    \draw[fill=gray!20, thick] (0.8,0) ellipse (0.45 and 0.25);
    \draw[thick] (-0.35,0) -- (0.35,0);
  \end{scope}
  \begin{scope}[shift={(5.0,1.4)}, scale=0.28, rotate=-20]
    \draw[fill=gray!20, thick] (-1.0,0) ellipse (0.55 and 0.18);
    \draw[fill=gray!20, thick] (1.0,0) ellipse (0.55 and 0.18);
    \draw[thick] (-0.45,0) -- (0.45,0);
  \end{scope}
  \begin{scope}[shift={(5.35,3.05)}, scale=0.28, rotate=45]
    \draw[fill=gray!20, thick] (-0.8,0) ellipse (0.45 and 0.25);
    \draw[fill=gray!20, thick] (0.8,0) ellipse (0.45 and 0.25);
    \draw[thick] (-0.35,0) -- (0.35,0);
    \draw[red!80!black, thick, ->] (0.0,0.65) arc (90:300:0.45);
  \end{scope}
  \draw[red, thick, ->] (4.7,2.55) -- (5.15,2.85);
  \node[align=center] at (3.2,0.65) {closed shape loop\\can have nonzero holonomy};
  \node[red!80!black, align=center] at (5.6,2.45) {net rigid\\rotation};
\end{tikzpicture}
\caption{Gauge kinematics: a closed loop in shape space can produce a net rigid
motion. The start and end shape can be the same while the rigid placement has
changed; the effect is the holonomy of a connection.}
\label{fig:gauge-shape-cycle-holonomy}
\end{figure}

\subsection{Mechanical Connection From Zero Momentum}
\label{subsec:mechanical-connection-zero-momentum}

The mechanical connection is obtained by solving a constraint, not by adding a
new force law. If the body has no total rigid momentum, then any internal shape
velocity that would otherwise create angular or linear momentum must be
cancelled by an accompanying rigid velocity. The connection is the linear rule
that performs this cancellation at each shape.

Let the body velocity in the moving frame be
\[
  \xi=g^{-1}\dot g\in\mathfrak g,
\]
where \(\mathfrak g\) is the Lie algebra of \(G\). Suppose the kinetic energy
has the block form
\[
  T
  =
  \frac12
  \begin{pmatrix}
  \xi\\ \dot s
  \end{pmatrix}^T
  \begin{pmatrix}
  M(s) & A(s)\\
  A(s)^T & m(s)
  \end{pmatrix}
  \begin{pmatrix}
  \xi\\ \dot s
  \end{pmatrix}.
\]
The momentum conjugate to rigid motion is
\[
  \Pi=\frac{\partial T}{\partial \xi}=M(s)\xi+A(s)\dot s.
\]
The symbol \(\Pi\) is the total rigid-motion momentum written in body
variables: for \(G=\SE(3)\) it packages total linear and angular momentum.
If the total angular and linear momentum vanish,
\[
  \Pi=0,
\]
as they do for an isolated body prepared with zero net external impulse, then
\(\Pi\) remains zero by conservation and
\[
  \xi=-M(s)^{-1}A(s)\dot s.
\]
Writing \(\dot s=\dot s^\alpha\partial_\alpha\), define the connection
\[
  \boxed{
  \mathcal A_\alpha(s)=M(s)^{-1}A_\alpha(s).
  }
\]
The rigid velocity is then
\[
  \boxed{
  g^{-1}\dot g=-\mathcal A_\alpha(s)\dot s^\alpha.
  }
\]
This is the cleanest derivation of the gauge field: it is the rule that tells a
zero-momentum deforming body how much rigid motion is forced by an infinitesimal
shape change.

The minus sign is part of the reconstruction convention. A positive connection
component means that the chosen shape change would create rigid momentum in the
body frame; the actual zero-momentum motion supplies the opposite rigid
velocity to cancel it. This is the mechanical meaning behind the gauge-theory
formula.

\subsection{Explicit Inertia-Operator Form}
\label{subsec:explicit-gauge-equation}

A concrete \(\SO(3)\) version of the same construction is useful because it
turns the abstract mechanical connection into a linear algebra problem. Let
\[
  A=R^{-1}\dot R\in\mathfrak{so}(3)
\]
be the angular velocity matrix of the body frame. Let \(Q=Q^T\) encode the
instantaneous shape through the second moment tensor
\[
  Q_{ij}=\int_B\rho(X)X_iX_j\,\dd^3X
\]
in the chosen internal frame. The corresponding inertia tensor is
\[
  \widetilde I_{ij}
  =
  \delta_{ij}\sum_{k=1}^3Q_{kk}-Q_{ij}.
\]
Let \(Y(X,s(t))\) be the position of a material point in the chosen internal
frame, with the center of mass at the origin. The spatial position is
\(x=RY\). In the body frame,
\[
  R^{-1}\dot x=A Y+\dot Y.
\]
The first term is the rigid velocity required by the frame motion; the second
term is internal deformation. The body angular momentum is therefore
\[
  M_{\mathrm{tot}}
  =
  \int_B\rho\,Y\times(AY+\dot Y)\,\dd^3X.
\]
With the hat convention used for the rigid body,
\(\widehat v\,w=v\times w\), or
\[
  (\widehat v)_{ij}=-\epsilon_{ijk}v_k,
\]
the rigid part obeys the matrix identity
\[
  \widehat{\widetilde I\Omega}=AQ+QA,
  \qquad A=\widehat\Omega.
\]
The internal part is
\[
  J_{\mathrm{int}}
  =
  \int_B\rho\,Y\times\dot Y\,\dd^3X.
\]
For zero total angular momentum,
\[
  \widetilde I\Omega+J_{\mathrm{int}}=0.
\]
It is convenient to define the deformation-generated momentum that the rigid
motion must cancel by
\[
  \widetilde J=-J_{\mathrm{int}}.
\]
Then the zero-momentum condition takes the linear matrix form
\[
  \boxed{
  AQ+QA=\widehat{\widetilde J}.
  }
\]
Solving this equation is the explicit finite-dimensional gauge connection:
the angular velocity \(A\) is determined by the shape velocity through
\(\widetilde J\).

\begin{derivation}[Solving the finite-dimensional gauge equation]
Write
\[
  A=\widehat\Omega,\qquad A_{ij}=-\epsilon_{ijk}\Omega_k.
\]
To verify the identity \(AQ+QA=\widehat{\widetilde I\Omega}\), diagonalize the
symmetric tensor \(Q\) at the instant under consideration:
\[
  Q=\diag(q_1,q_2,q_3).
\]
Then
\[
  (AQ+QA)_{12}=-(q_1+q_2)\Omega_3,\quad
  (AQ+QA)_{23}=-(q_2+q_3)\Omega_1,\quad
  (AQ+QA)_{31}=-(q_3+q_1)\Omega_2.
\]
But the inertia tensor has principal moments
\[
  \widetilde I_1=q_2+q_3,\qquad
  \widetilde I_2=q_3+q_1,\qquad
  \widetilde I_3=q_1+q_2,
\]
so these entries are exactly those of
\(\widehat{\widetilde I\Omega}\). Since both sides transform tensorially under
orthogonal changes of the internal frame, the identity holds in any frame.
At zero total angular momentum, the rigid angular momentum must equal
\(\widetilde J\), the negative of the internal contribution. The linear map
from angular velocity \(\Omega\) to rigid angular momentum is precisely the
inertia tensor:
\[
  \widetilde J_i=\widetilde I_{ij}\Omega_j.
\]
Equivalently,
\[
  \Omega_i=(\widetilde I^{-1})_{ij}\widetilde J_j.
\]
Therefore
\[
  \boxed{
  A_{ij}
  =
  -\epsilon_{ijk}
  (\widetilde I^{-1})_{k\ell}\widetilde J_\ell.
  }
\]
This is the concrete \(\SO(3)\) deforming-body gauge formula. In the abstract
notation above, \(M(s)\) is the inertia operator
\(\widetilde I(s)\), while the cross term \(A(s)\dot s\) is the internally
generated angular momentum \(\widetilde J(s,\dot s)\). The helicopter rotor
laboratory in Section~\ref{sec:lab-helicopter-gauge} is the finite-dimensional
special case where the shape variables are two rotor angles and this abstract
inertia-operator equation reduces to a pair of scalar connection components.
\end{derivation}

\subsection{Gauge Transformation And Holonomy}
\label{subsec:gauge-holonomy}

Under a shape-dependent change of body frame \(h(s)\), the connection transforms
as
\[
  \boxed{
  \mathcal A'
  =
  \Ad_{h^{-1}}\mathcal A-h^{-1}\dd h.
  }
\]
The sign of the inhomogeneous term follows from the reconstruction convention
\(g^{-1}\dot g=-\mathcal A(\dot s)\). With the opposite convention, the same
geometric connection is written with the more familiar plus sign. The base space
is shape space and the structure group is the group of rigid motions.

For a path \(s(t)\), the rigid motion is obtained by integrating
\[
  g^{-1}\dot g=-\mathcal A_\alpha(s)\dot s^\alpha.
\]
For a closed shape cycle \(\Gamma\), the net displacement is
\[
  g(T)
  =
  g(0)\,
  \mathcal P\exp\left(
    -\oint_\Gamma \mathcal A
  \right),
\]
where \(\mathcal P\) denotes path ordering. The result can be nontrivial even
though the shape returns to its starting point.

The curvature of the connection is
\[
  \boxed{
  \mathcal F
  =
  \dd\mathcal A+\mathcal A\wedge\mathcal A,
  }
\]
or in components,
\[
  \mathcal F_{\alpha\beta}
  =
  \partial_\alpha\mathcal A_\beta
  -
  \partial_\beta\mathcal A_\alpha
  +
  [\mathcal A_\alpha,\mathcal A_\beta].
\]
For a small loop enclosing area element
\(\Delta S^{\alpha\beta}\),
\[
  \Delta g
  \approx
  \exp\left(
    -\mathcal F_{\alpha\beta}\Delta S^{\alpha\beta}
  \right).
\]

\begin{figure}[htbp]
\centering
\begin{tikzpicture}[scale=1.0, >=Latex]
  \draw[->] (-0.2,0) -- (6.3,0) node[right] {\(s^1\)};
  \draw[->] (0,-0.2) -- (0,3.7) node[above] {\(s^2\)};
  \coordinate (a) at (1.6,1.0);
  \coordinate (b) at (4.4,1.0);
  \coordinate (c) at (4.4,2.7);
  \coordinate (d) at (1.6,2.7);
  \fill[blue!10] (a) -- (b) -- (c) -- (d) -- cycle;
  \draw[blue!70!black, thick, ->] (a) -- node[below] {\(\delta s^1\)} (b);
  \draw[blue!70!black, thick, ->] (b) -- node[right] {\(\delta s^2\)} (c);
  \draw[blue!70!black, thick, ->] (c) -- (d);
  \draw[blue!70!black, thick, ->] (d) -- (a);
  \draw[red!75!black, thick, ->] (3.0,1.85) arc (210:-80:0.45);
  \node[red!75!black, align=center] at (5.15,2.35)
    {net rigid motion\\\(-\mathcal F_{12}\delta s^1\delta s^2\)};
  \node[align=center] at (3.0,0.45)
    {same infinitesimal shape loop,\\different final placement};
\end{tikzpicture}
\caption{Curvature is the local density of holonomy. Traversing a small
rectangle in shape space returns to the same shape but can leave a residual
rigid displacement proportional to the curvature through the enclosed area.}
\label{fig:gauge-curvature-small-loop}
\end{figure}

Thus curvature is the local density of geometric phase. This is the same
mathematical idea as Hannay angle, Berry phase, and the rigid-body falling-cat
effect, now written for deforming bodies.

Flat connections can still produce frame changes along open paths, but a small
closed loop detects curvature. This is why the curvature, rather than the
connection coefficients themselves, is the local observable of the gauge
description. Different choices of body frame change \(\mathcal A\), while the
holonomy around a loop changes by conjugation and therefore represents the same
physical displacement.

\section{Swimming At Low Reynolds Number}
\label{sec:low-reynolds-swimming}

Shapere and Wilczek's low-Reynolds-number swimming problem gives a particularly
physical gauge theory. At very small Reynolds number, inertia drops out of
Navier-Stokes and the surrounding fluid satisfies Stokes equations:
\[
  -\nabla p+\mu\Delta u=0,\qquad \nabla\cdot u=0.
\]
Boundary velocity is prescribed by the swimmer's rigid motion plus deformation
rate. Because the Stokes problem is linear, the force and torque on the body
depend linearly on:
\[
  \xi=g^{-1}\dot g
  \qquad\text{and}\qquad
  \dot s.
\]
A freely swimming body has zero net force and torque, so the same algebra as
above gives
\[
  \xi=-\mathcal A_\alpha(s)\dot s^\alpha.
\]
The connection is now computed from a hydrodynamic boundary-value problem
rather than from an inertia tensor. A reciprocal one-parameter stroke encloses
no area in shape space and produces no net motion: this is Purcell's scallop
theorem in gauge language. A two-parameter stroke can enclose area and swim by
curvature.

\section{Local Frames In Elastic Continua}
\label{sec:local-frames-elastic-continua}

The finite-dimensional gauge picture becomes a field theory when every material
point carries a local frame. A Cosserat body has, in addition to
\(\varphi(X,t)\), a rotation field
\[
  R(X,t)\in\SO(3).
\]
The directors are
\[
  d_a(X,t)=R(X,t)E_a,
\]
where \(E_a\) is a fixed reference basis. The rotational connection is
\[
  A_A=R^{-1}\partial_A R\in\mathfrak{so}(3),
  \qquad
  A_t=R^{-1}\partial_t R.
\]
The translational strain measured in the local frame is
\[
  \Gamma_A=R^{-1}\partial_A\varphi.
\]
For a rigid body, \(R\) depends only on time. For a Cosserat continuum, \(R\)
depends on \(X\), so neighboring material points can rotate relative to one
another.

This is the field-theoretic version of the rigid-body frame issue. In the
rigid body there is one moving frame; in a Cosserat body there is a moving frame
at every material point. The connection \(A_A\) compares nearby frames, while
\(A_t\) records how a given local frame evolves in time.

\begin{definition}[Local frame gauge transformation]
Let \(G(X,t)\in\SO(3)\) be a change of local material frame:
\[
  R'(X,t)=R(X,t)G(X,t).
\]
Then
\[
  A'_A=G^{-1}A_AG+G^{-1}\partial_A G,
  \qquad
  \Gamma'_A=G^{-1}\Gamma_A.
\]
Thus \(A_A\) is a gauge connection and \(\Gamma_A\) is a vector-valued
matter field in the local frame.
\end{definition}

\section{Curvature, Torsion, And Defects}
\label{sec:curvature-torsion-defects}

If a connection is globally of the form
\[
  A_A=R^{-1}\partial_A R,
\]
then it is pure gauge and satisfies the Maurer-Cartan identity:
\[
  F_{AB}
  =
  \partial_AA_B-\partial_BA_A+[A_A,A_B]
  =
  0.
\]
Nonzero curvature means that local frames cannot be integrated into a single
smooth orientation field without singularities or defects. In defect theory,
curvature is associated with disclinations, or rotational defects.

To describe translational defects, introduce a local coframe \(\theta^a\), which
measures material line elements in a local frame. The connection one-forms
\(\omega^a{}_b\) are the components of the same \(\mathfrak{so}(3)\)-valued
connection written above: if the matrix entries of \(A_A\) are
\((A_A)^a{}_b\), then
\[
  \omega^a{}_b=(A_A)^a{}_b\,\dd X^A.
\]
With this notation, the torsion two-form is
\[
  T^a=\dd\theta^a+\omega^a{}_b\wedge\theta^b.
\]
Torsion measures closure failure of infinitesimal parallelograms and is
associated with dislocation density. In this language:
\[
\begin{aligned}
  \text{curvature} &\longleftrightarrow \text{rotational mismatch},\\
  \text{torsion} &\longleftrightarrow \text{translational mismatch},\\
  \text{incompatible strain} &\longleftrightarrow \text{nonintegrable local data}.
\end{aligned}
\]
The gauge theory is not decorative terminology. It records which descriptions
depend on arbitrary local choices of frame and which quantities are invariant.

\section{Cosserat Rod As A One-Dimensional Gauge Theory}
\label{sec:cosserat-rod}

A planar rod has material coordinate \(s\), centerline \(r(s)\in\R^2\), and
orientation angle \(\theta(s)\). Its director is
\[
  d_1=(\cos\theta,\sin\theta).
\]
For an inextensible unshearable rod,
\[
  r'(s)=d_1(s).
\]
The connection is the scalar curvature
\[
  \kappa(s)=\theta'(s).
\]
Given \(\kappa\),
\[
  \theta(s)=\theta(0)+\int_0^s\kappa(\sigma)\,\dd\sigma,
\]
and
\[
  r(s)=r(0)+\int_0^s
  (\cos\theta(\sigma),\sin\theta(\sigma))\,\dd\sigma.
\]
This is the reconstruction problem: a gauge field \(\kappa\,\dd s\) determines
shape after choosing initial position and frame.

The bending energy
\[
  E_{\mathrm{bend}}
  =
  \frac12\int_0^L B(\kappa-\kappa_0)^2\,\dd s
\]
compares actual curvature with preferred curvature. Here \(B\) is bending
stiffness and \(\kappa_0\) is intrinsic curvature. In gauge language, the
energy penalizes mismatch between a connection and a preferred connection.

\section{Conceptual Summary}
\label{sec:elastic-conceptual-summary}

The progression from rigid body to elastic gauge theory is:
\[
\begin{aligned}
  \text{rigid body:}&\quad R(t),\quad R^{-1}\dot R,\\
  \text{deformable body:}&\quad g(t),\ s(t),\quad g^{-1}\dot g=-\mathcal A(s)\dot s,\\
  \text{Cosserat continuum:}&\quad R(X,t),\quad R^{-1}\partial_A R,\\
  \text{defect theory:}&\quad A,\quad F=\dd A+A\wedge A,\quad T=\dd\theta+\omega\wedge\theta.
\end{aligned}
\]
The first line is ordinary rigid-body mechanics. The second is
Shapere-Wilczek gauge kinematics over shape space. The third attaches a local
rigid body to every material point. The fourth allows the local frame and
translation data to be incompatible, producing curvature and torsion as physical
defect densities.

\section{Associated Demos}
\label{sec:elastic-associated-demos}

\begin{itemize}
\item \path{demos/python/cosserat_rod_demo.py}, a planar frame reconstruction
from curvature.
\item \path{demos/mathematica/GaugeDeformingBody.wl}, a finite-dimensional
gauge-connection toy model.
\item Suggested extension: a two-shape-parameter swimmer showing geometric phase
as the area integral of connection curvature.
\end{itemize}
